\section{Decision Support}\label{sec:decision-support}

\subsection{AI-based}

\subsection{Predictive Maintenance}

\subsection{Explainable AI}

In \cite{alexander-an-interrogative-survey-of-explainable-ai-in-manufacturing}, the authors define interpretable and explainable \ac{ai} and review their use cases in manufacturing.
The authors describe interpretable \ac{ai} as methods that are inherently easy to understand due to their data-processing architecture.
For instance, random forest methods fall into this category, as it is relatively straightforward to trace a decision through the rules of each tree, albeit a cumbersome task when dealing with many trees.
In contrast, explainable \ac{ai} refers to techniques used to interpret decisions made by complex models such as \ac{dl} models.
These techniques often rely on local explanations of individual model outputs, such as \ac{shap} values, Grad-CAM, or \ac{lime}.

\subsection{Identified Areas Where AI can be Used}

\subsubsection{Before Manufacturing}
Product design verification (can then Product be manufactured?).
Production planning and optimization (plan the execution in an optimal way).
Demand forecasting (how many to produce), raw materials and components forecasting (how much material do we need?).

\subsubsection{During Manufacturing}
Keep production lines running.
Defect free production.
Machine health monitoring (predictive maintenance) (maintenance costs a lot).
Fault diagnostics (why).

\subsubsection{After Manufacturing}
QA (check for defects).
Testing.

\subsection{Challenges in AI}

\begin{itemize}
  \item Trusting and understanding black box models.
\end{itemize}
